% Avsnitt 6.8, ur lectures/L06/README.md.

\section{Sammanfattning}
\label{c6:sec:review}

\needspace{6\baselineskip}
\subsection*{Det här ska du kunna nu}
\begin{itemize}
  \item Bygga en 4-bitars räknare som grindnät, och förklara, med fingret på din egen ritning,
    varför ingenting i den får räkningen att återgå till noll.
  \item Implementera samma räknare i VHDL, förklara överslaget som en konsekvens av registrets
    bitbredd, och säga varför en simulator invänder mot ett överslag som hårdvaran utför utan att
    blinka.
  \item Implementera ett skiftregister i både SIPO- och PISO-utförande, veta vilket som passar
    vilket jobb, och veta vilket håll boken skiftar åt och varför det spelar roll att hålla sig
    till ett.
  \item Kombinera en räknare med ett skiftregister för att upptäcka att ett komplett ord anlänt,
    och verifiera en klockad design genom att köra dess testbänk.
\end{itemize}

\needspace{6\baselineskip}
\subsection*{Frågor att testa dig själv med}
\begin{enumerate}
  \item Varför återgår en räknare av typen \code{natural range 0 to 15} till \code{0} efter
    \code{15} utan någon uttrycklig ”om maxvärdet, nollställ”-logik? Vad måste gälla för
    intervallet? Vilket element utför överslaget?
  \item En VHDL-simulator avbryter på det överslag din CircuitVerse-ritning utför utan invändning.
    Vilken av de två modellerar hårdvaran, och vad säger det om att lita på någondera ensam?
  \item Vad är den praktiska skillnaden mellan ett SIPO- och ett PISO-skiftregister? Vilket skulle
    du gripa efter för att ta emot en ström av seriella bitar, och vilket för att driva en kedja av
    lysdioder från ett fast mönster?
  \item Ett skiftregister och en räknare är båda ”N vippor som delar en klocka”. Det som skiljer är
    vad som matar varje vippas ingång. Beskriv den skillnaden för respektive.
  \item Varför är \code{data_ready} i \code{serial_rx8} en puls på en klockcykel snarare än en nivå
    som blir liggande hög när en byte anlänt?
  \item Varför måste \code{shift_enable} grinda biträknaren och inte bara skiftningen?
  \item Kortdemonstrationen driver \code{shift_enable} från en knapp i stället för att sakta ned
    klockan. Varför är nedsaktad klocka fel svar, och vad skulle gå sönder om knappen kopplades
    till \code{shift_enable} utan att först passera \code{button_sync}?
  \item \code{data_ready} behöver en egen vippa innan en lysdiod kan visa den. Vilket kapitels
    problem är det, i annan skepnad?
\end{enumerate}

\needspace{6\baselineskip}
\subsection*{Blicken framåt}
\Chapref{c7:ch} går vidare med:
\begin{itemize}
  \item Timers: en räknare med ett målvärde, vilket är det som gör klockcykler till riktig tid.
  \item Byggd genom att lägga exakt två saker till räknaren du ritade här, en jämförare och en
    nollställning, så spara det CircuitVerse-projektet.
  \item Den vandrande lysdioden: det här kapitlets skiftregister, taktat av den timern och startat
    av en synkroniserad knapp.
\end{itemize}
